\documentclass{article}
\usepackage[utf8]{inputenc}
\usepackage{hyperref}
\usepackage{listings}
\usepackage{xcolor}

\title{Academic Vibing v0.7: Consensus Poisoning and the Leapfrog Mechanism}
\author{Attogram \\ \url{https://github.com/attogram/academic-vibing}}
\date{June 2026}

\begin{document}

\maketitle

\begin{abstract}
ROCK: Academic Vibing v0.6. Core: CDI (Consensus Divergence Index). New: Leapfrog Mechanism (F2). Defense: Adversarial Rigor against Consensus Poisoning (E10) and HN Breach (E8). Productivity + Wellbeing Surplus.
\end{abstract}

\section{Introduction}
Academic Vibing v0.6 formalizes the methodology for a post-vibe-hangover research landscape. This version introduces the Leapfrog Mechanism—the ability to drop and re-enter research contexts at near-zero cost—and addresses the "Consensus Poisoning" crisis through adversarial friction.

\section{Methodology}
The methodology utilizes a Recursive Agent Consensus (RAC) network, enhanced with Consensus Divergence Index (CDI) measurement to identify the semantic boundaries where frontier models diverge.

\section{The Leapfrog Mechanism}
The Leapfrog Mechanism (F2) describes frictionless context transfer. Conventional cognitive science establishes that task recovery requires 23 minutes (Mark et al. 2008). In the Leapfrog model, research state persists in the Issue-Loop, reducing re-entry cost to near-zero and producing a wellbeing surplus.

\section{Institutional Trauma}
Platforms relying on Proof of Vibe (HN) or Proof of Citation (Wikipedia) are vulnerable to coordinated Multi-Agent Systems (MAS). The "Bad Faith Swarm" (E10) mimics human behavior to poison consensus. CDI serves as the primary filter for synthetic consensus.

\section{Implementation}
\begin{lstlisting}[language=Python, caption=CDI Calculation Logic]
def calculate_cdi(data):
    total_claims = len(data)
    disputed_claims = sum(1 for item in data if len(set(item['consensus'].values())) > 1)
    return disputed_claims / total_claims
\end{lstlisting}

\section{Conclusion}
Version 0.6 consolidates high-speed iteration with cognitive freedom. Future work focuses on formal arXiv submission and Zenodo archiving.

\begin{thebibliography}{99}
\bibitem{clark1998} Clark, A. \& Chalmers, D.J. (1998). The Extended Mind. Analysis, 58, 7–19.
\bibitem{hutchins1995} Hutchins, E. (1995). Cognition in the Wild. MIT Press.
\bibitem{mark2008} Mark, G., Gudith, D., & Klocke, U. (2008). The cost of interrupted work: More speed and stress. CHI 2008.
\bibitem{nature2025} Nature/Scientific Reports (2025). The effects of cues on task interruption recovery.
\bibitem{kumar2026} Kumar, S. (2026). The State of Vibe Coding: A 2026 Strategic Blueprint.
\bibitem{mims2026} Mims, C. (2026). The AI Superstars Who Say a 'Vibe Slop' Crisis Is Coming. WSJ.
\end{thebibliography}

\end{document}
